\documentclass[10pt,a4paper]{amsart}
\usepackage{amsmath,amssymb,amsthm,graphicx}
% For commutative diagrams we use Xy-pic
\usepackage[all]{xy}

% We use 2cell for 2-commutative diagrams.
\xyoption{2cell}
% Theorem environments.
%
\theoremstyle{plain}
\newtheorem{theorem}[subsection]{Theorem}
\newtheorem{proposition}[subsection]{Proposition}
\newtheorem{lemma}[subsection]{Lemma}

\theoremstyle{definition}
\newtheorem{definition}[subsection]{Definition}
\newtheorem{example}[subsection]{Example}
\newtheorem{exercise}[subsection]{Exercise}
\newtheorem{situation}[subsection]{Situation}

\theoremstyle{remark}
\newtheorem{remark}[subsection]{Remark}
\newtheorem{remarks}[subsection]{Remarks}

\numberwithin{equation}{subsection}

% Macros
%
\def\lim{\mathop{\mathrm{lim}}\nolimits}
\def\colim{\mathop{\mathrm{colim}}\nolimits}
\def\Spec{\mathop{\mathrm{Spec}}}
\def\Hom{\mathop{\mathrm{Hom}}\nolimits}
\def\Ext{\mathop{\mathrm{Ext}}\nolimits}
\def\SheafHom{\mathop{\mathcal{H}\!\mathit{om}}\nolimits}
\def\SheafExt{\mathop{\mathcal{E}\!\mathit{xt}}\nolimits}
\def\Sch{\mathit{Sch}}
\def\Mor{\mathop{\mathrm{Mor}}\nolimits}
\def\Ob{\mathop{\mathrm{Ob}}\nolimits}
\def\Sh{\mathop{\mathit{Sh}}\nolimits}
\def\NL{\mathop{N\!L}\nolimits}
\def\CH{\mathop{\mathrm{CH}}\nolimits}
\def\proetale{{pro\text{-}\acute{e}tale}}
\def\etale{{\acute{e}tale}}
\def\QCoh{\mathit{QCoh}}
\def\Ker{\mathop{\mathrm{Ker}}}
\def\Im{\mathop{\mathrm{Im}}}
\def\Coker{\mathop{\mathrm{Coker}}}
\def\Coim{\mathop{\mathrm{Coim}}}

% Boxtimes
%
\DeclareMathSymbol{\boxtimes}{\mathbin}{AMSa}{"02}

%
% Macros for moduli stacks/spaces
%
\def\QCohstack{\mathcal{QC}\!\mathit{oh}}
\def\Cohstack{\mathcal{C}\!\mathit{oh}}
\def\Spacesstack{\mathcal{S}\!\mathit{paces}}
\def\Quotfunctor{\mathrm{Quot}}
\def\Hilbfunctor{\mathrm{Hilb}}
\def\Curvesstack{\mathcal{C}\!\mathit{urves}}
\def\Polarizedstack{\mathcal{P}\!\mathit{olarized}}
\def\Complexesstack{\mathcal{C}\!\mathit{omplexes}}
% \Pic is the operator that assigns to X its picard group, usage \Pic(X)
% \Picardstack_{X/B} denotes the Picard stack of X over B
% \Picardfunctor_{X/B} denotes the Picard functor of X over B
\def\Pic{\mathop{\mathrm{Pic}}\nolimits}
\def\Picardstack{\mathcal{P}\!\mathit{ic}}
\def\Picardfunctor{\mathrm{Pic}}
\def\Deformationcategory{\mathcal{D}\!\mathit{ef}}
\usepackage{cleveref}
\crefname{originallemma}{original lemma}{original lemmas}
\Crefname{originallemma}{Original lemma}{Original lemmas}
\makeatletter
\newlabel{lemma-prepare-tame-symbol}{{0EAG}{source}{}{sourcepointer-1}{}}
\newlabel{lemma-prepare-tame-symbol@cref}{{[originallemma][0][]0EAG}{[1][0][]source}{}{}{}}
\makeatother
\crefname{originallemma}{original lemma}{original lemmas}
\Crefname{originallemma}{Original lemma}{Original lemmas}
\makeatletter
\newlabel{lemma-glue-at-max}{{0EAE}{source}{}{sourcepointer-2}{}}
\newlabel{lemma-glue-at-max@cref}{{[originallemma][0][]0EAE}{[1][0][]source}{}{}{}}
\makeatother
\crefname{originallemma}{original lemma}{original lemmas}
\Crefname{originallemma}{Original lemma}{Original lemmas}
\makeatletter
\newlabel{fields-lemma-trace-and-norm-tower}{{0BIJ}{source}{}{sourcepointer-3}{}}
\newlabel{fields-lemma-trace-and-norm-tower@cref}{{[originallemma][0][]0BIJ}{[1][0][]source}{}{}{}}
\makeatother
\title{The tame symbol on a one-dimensional local ring is the norm of the tame symbols of its finite local extensions}
\author{The Stacks Project Authors}
\date{stacks-project:chow.tex:local-snapshot-20261011; GNU FDL 1.2 or later}
\begin{document}
\maketitle
\noindent\textit{Editorial note.} All delivered mathematical text is the author's own, copied from the pinned Stacks Project chapter file chow.tex; no statement or proof step was rewritten, reordered or removed.\par\medskip
\noindent\textit{Editorial note.} The preamble is the chapter's own native preamble excerpt used for this source: the Xy-pic setup at lines 14-18 and the theorem declarations with the mathematical macros at lines 154-217 of sources/stacks-chow/preamble.tex. The quoted class line and the xr-hyper externaldocument block are not included because the assembler supplies a plain amsart wrapper; no macro was invented or redefined.\par\medskip
\noindent\textit{Editorial note.} The two same-chapter references are located at their \textbackslash{}label lines inside the pinned chapter file; the cross-chapter pointer is located in the project's own tag list staged in the snapshot (sources/stacks-chow/tags.txt), at the verbatim line carrying its canonical tag.\par\medskip
\noindent\textit{Allowed standard prerequisites (editorial).}\begin{itemize}
\item Preparation of the local extensions so that both tame symbols become computable. Same chapter, Lemma labelled lemma-prepare-tame-symbol (source line 642, Stacks tag 0EAG, already counted by the card chow-prepare-tame-symbol); invoked at proof lines 967-968. Supplies the extensions B\_j contained in C\_j in which a and b admit local factorisations of the displayed shape; the computation of the two norms on top of it is retained in the counted span.
\item Gluing finitely many finite local extensions into one finite extension. Same chapter, Lemma labelled lemma-glue-at-max (source line 526, Stacks tag 0EAE, counted earlier in this night's work); invoked at proof lines 969-970. Supplies the single finite extension B contained in C with C\_j isomorphic to C\_\{m\_j\}, which is what makes the product over the maximal ideals of C meaningful.
\item Transitivity of the norm in a tower of field extensions. Fields, Lemma labelled fields-lemma-trace-and-norm-tower (canonical Stacks tag 0BIJ); invoked at proof lines 994-996. Supplies the passage from the double product over (j, k) to the product over j of norms for kappa(m\_\{j,k\})/kappa(m\_j), which is the last step.
\end{itemize}\medskip
\ifdefined\hypertarget\hypertarget{sourcepointer-1}{}\fi
\noindent\textit{Source reference.} Original lemma 0EAG: sources/stacks-chow/pinned-original-source.tex, source lines 642--642. Same-chapter pointer invoked at proof lines 967-968 to choose the extensions in which a and b admit local factorisations; the norm computation built on it is retained verbatim in the counted span.\par\medskip
\ifdefined\hypertarget\hypertarget{sourcepointer-2}{}\fi
\noindent\textit{Source reference.} Original lemma 0EAE: sources/stacks-chow/pinned-original-source.tex, source lines 526--526. Same-chapter pointer invoked at proof lines 969-970 to glue the prepared local extensions into one finite extension; the regrouping of the two products is retained in the counted span.\par\medskip
\ifdefined\hypertarget\hypertarget{sourcepointer-3}{}\fi
\noindent\textit{Source reference.} Original lemma 0BIJ: sources/stacks-chow/tags.txt, source lines 14492--14492. Cross-chapter standard prerequisite invoked once at proof lines 994-996 for transitivity of the norm; the two displays it compares are retained in the counted span.\par\medskip
\begin{lemma}
\label{lemma-norm-down-tame-symbol}
Let $(A, \mathfrak m)$ be a Noetherian local ring of dimension $1$.
Let $A \subset B$ be a finite ring extension with $B/A$
annihilated by a power of $\mathfrak m$ and $\mathfrak m$ not
an associated prime of $B$.
For $a, b \in A$ nonzerodivisors we have
$$
\partial_A(a, b) = \prod
\text{Norm}_{\kappa(\mathfrak m_j)/\kappa(\mathfrak m)}(\partial_{B_j}(a, b))
$$
where the product is over the maximal ideals $\mathfrak m_j$ of $B$
and $B_j = B_{\mathfrak m_j}$.
\end{lemma}

\begin{proof}
Choose $B_j \subset C_j$ as in
Lemma \ref{lemma-prepare-tame-symbol} for $a, b$.
By Lemma \ref{lemma-glue-at-max} we can choose a finite ring
extension $B \subset C$ with $C_j \cong C_{\mathfrak m_j}$ for all $j$.
Let $\mathfrak m_{j, k} \subset C$ be the maximal ideals of $C$
lying over $\mathfrak m_j$. Let
$$
a = u_{j, k}\pi_{j, k}^{f_{j, k}},\quad
b = v_{j, k}\pi_{j, k}^{g_{j, k}}
$$
be the local factorizations which exist by our choice of
$C_j \cong C_{\mathfrak m_j}$. By definition we have
$$
\partial_A(a, b) = 
\prod\nolimits_{j, k}
\text{Norm}_{\kappa(\mathfrak m_{j, k})/\kappa(\mathfrak m)}
((-1)^{f_{j, k}g_{j, k}}u_{j, k}^{g_{j, k}}v_{j, k}^{-f_{j, k}}
\bmod \mathfrak m_{j, k})^{m_{j, k}}
$$
and
$$
\partial_{B_j}(a, b) = 
\prod\nolimits_k
\text{Norm}_{\kappa(\mathfrak m_{j, k})/\kappa(\mathfrak m_j)}
((-1)^{f_{j, k}g_{j, k}}u_{j, k}^{g_{j, k}}v_{j, k}^{-f_{j, k}}
\bmod \mathfrak m_{j, k})^{m_{j, k}}
$$
The result follows by transitivity of norms
for $\kappa(\mathfrak m_{j, k})/\kappa(\mathfrak m_j)/\kappa(\mathfrak m)$, see
Fields, Lemma \ref{fields-lemma-trace-and-norm-tower}.
\end{proof}

\end{document}
